\section*{Overview}

Monotonic stacks and monotonic queues are the ``throw away dominated candidates immediately'' tools. They show up in
array problems where the answer depends on the nearest larger/smaller element or on the minimum/maximum inside a window.

\section*{When to Use It}

Use them when:

\begin{itemize}
  \item candidates arrive in a fixed order,
  \item once one candidate dominates another, the dominated one will never be useful again,
  \item the query asks for a nearest blocker or for the best element in a moving window.
\end{itemize}

\section*{Core Idea}

Keep indices in a stack or deque whose corresponding values stay monotone.

\begin{itemize}
  \item For nearest greater/smaller problems, the structure is usually a stack.
  \item For sliding windows, it is usually a deque because indices leave from the front and dominated indices leave from
  the back.
\end{itemize}

\section*{Key Insight}

Each index enters once and leaves once. The linear-time guarantee is not accidental; it comes from the dominance rule.
If an index is popped, it is gone forever.

\section*{Worked Problem}

\paragraph{Problem 1.}
For every position \(i\), find the next position to the right whose value is strictly smaller.

\paragraph{Problem 2.}
For every window of length \(k\), report its minimum.

\paragraph{Why monotonic structures fit.}
In both problems, worse candidates can be discarded permanently:

\begin{itemize}
  \item a larger element cannot be the next smaller element for anyone once a smaller element appears before it,
  \item an older larger element can never be the minimum of a future window if a newer smaller element is already inside.
\end{itemize}

\section*{Correctness Intuition}

The invariant is the whole proof:

\begin{itemize}
  \item stack values stay monotone, so the first surviving candidate in the needed direction is the nearest valid one,
  \item deque values stay monotone and indices stay increasing, so the front is always the best valid candidate inside
  the current window.
\end{itemize}

\section*{Complexity Analysis}

Both the stack and deque versions run in \(O(n)\) time and \(O(n)\) memory.

\section*{Implementation}

The code includes:

\begin{itemize}
  \item next strictly smaller element to the right,
  \item sliding-window minimum.
\end{itemize}

\section*{Common Pitfalls}

\begin{itemize}
  \item Choosing the wrong inequality and breaking ties incorrectly.
  \item Storing values instead of indices, then losing window-expiration information.
  \item Forgetting that some problems need strictly monotone order and others need non-strict order.
\end{itemize}

\section*{Variants / Extensions}

\begin{itemize}
  \item Largest rectangle in histogram.
  \item Monotone queue optimization in DP.
  \item Cartesian tree construction from the same stack discipline.
\end{itemize}

\section*{Practice Problems}

\begin{itemize}
  \item Next greater / next smaller element.
  \item Sliding-window minimum or maximum.
  \item Histogram and subarray-boundary problems driven by nearest blockers.
\end{itemize}

\section*{References}

\begin{itemize}
  \item \href{https://usaco.guide/gold/stacks?lang=cpp}{USACO Guide: Monotonic Stacks}.
  \item \href{https://codeforces.com/catalog}{Codeforces Catalog}.
  \item \href{https://usaco.guide/CPH.pdf}{Competitive Programmer's Handbook}.
\end{itemize}
